\documentclass[11pt]{article}
\usepackage[a4paper,margin=0.8in]{geometry}
\usepackage{amsmath,amssymb,amsthm,hyperref}
\hypersetup{hidelinks}
\newtheorem{theorem}{Theorem}
\newtheorem{lemma}{Lemma}
\title{A Gaussian envelope for Hermite polynomials}
\author{Research proof and Lean formalization}
\date{October 27, 2026}
\begin{document}
\maketitle
\begin{abstract}
We prove the Hermite-envelope conjecture of Guo and Liu. For odd degrees, the initial oscillator energy is slightly larger than the desired coefficient permits. A short sine comparison gives enough strict energy loss to remove this excess. A real-root moment inequality excludes nonzero stationary peaks after the conjectured cutoff and supplies the clamped tail. The proof uses classical oscillator, Wallis and polynomial tools; no new general machinery or priority claim is made.
\end{abstract}
\noindent\textbf{Verification status.} The literal source theorem is proved in Lean. Independent final review and the extracted submission-packet replay are pending; this candidate has discovery credit zero.

Let $H_n(x)=(-1)^n e^{x^2}(d/dx)^n e^{-x^2}$ be the physicists' Hermite polynomial, and set
\[
 E=2n+1,\qquad C_n=\frac{2^n\Gamma((n+1)/2)}{\sqrt\pi},\qquad
 a_n=\sqrt E\left(1-\frac\pi{2E}\right),\qquad
 g_n(t)=\left(1-\frac{t^2}{E}\right)^{-1/2}.
\]
The cutoff satisfies $0<a_n<\sqrt E$ for $n\geq1$. Guo--Liu's Conjecture 1, published as \cite[eq.~(29), p.~8]{GL}, is exactly the following statement.
\begin{theorem}\label{main}
For every integer $n\geq1$ and real $x>0$,
\[
 |H_n(x)e^{-x^2}|\leq C_n g_n(\min\{x,a_n\})e^{-x^2/2}.
\]
\end{theorem}

\section*{The bulk estimate}
Write $y(x)=H_n(x)e^{-x^2/2}$. The Hermite equation gives
\begin{equation}\label{energy}
 y''+(E-x^2)y=0,\quad B(x)=(E-x^2)y(x)^2+y'(x)^2,\quad B'(x)=-2xy(x)^2.
\end{equation}
For even $n$, $y'(0)=0$ and $|y(0)|=C_n$, so $B(x)\leq EC_n^2$ on $x\geq0$. This proves $|y(x)|\leq C_ng_n(x)$ for $0\leq x<\sqrt E$ in every even degree.

For $n=2m+1$, replace $y$ by $(-1)^m y$, so $y(0)=0$ and
$y'(0)=A=2(2m+1)(2m)!/m!>0$. Put $\gamma_m=A^2/(EC_n^2)$. The exact initial values and the Wallis limit give
\begin{align}
 \gamma_m&=\frac{\pi(2m+1)^2\binom{2m}{m}^2}{(4m+3)16^m},&
 \frac{\gamma_{m+1}}{\gamma_m}&=1-\frac1{4(m+1)^2(4m+7)},&\gamma_m&\longrightarrow1.\label{wallis}
\end{align}
For $\delta_m=\gamma_m(1-1/(2(4m+3)^2))$, direct cancellation yields
\[
 \delta_{m+1}-\delta_m=
 \gamma_m\frac{176m^2+440m+247}{8(m+1)^2(4m+3)^2(4m+7)^3}>0.
\]
Since $\delta_m\to1$, we have $\delta_m<1$. For $m\geq1$, also
$\gamma_m\leq\gamma_1=9\pi/28<99/98$. Thus
\begin{equation}\label{initial}
 A^2\left(1-\frac1{2E^2}\right)<EC_n^2,\qquad A^2\leq\frac{99}{98}EC_n^2\quad(m\geq1).
\end{equation}

Assume now $m\geq1$, so $E\geq7$. Put $x_0=\sqrt{2/E}$, $q_0=E-2/E$ and $c=\sqrt{1-2/E^2}$. On $0\leq x\leq x_0$,
\begin{equation}\label{sine}
 \frac A{\sqrt E}\sin(\sqrt E x)\leq y(x)\leq
 \frac A{\sqrt{q_0}}\sin(\sqrt{q_0}x).
\end{equation}
To justify positivity as well as comparison, $y'(0)>0$ first gives $y>0$ near zero. If its first positive zero $z$ were at most $x_0$, variation of constants would give
\[
 y(z)=\frac A{\sqrt E}\sin(\sqrt E z)+
 \int_0^z\frac{\sin(\sqrt E(z-t))}{\sqrt E}\,t^2y(t)\,dt>0,
\]
because $\sqrt E z\leq\sqrt2<\pi$. This is a contradiction. The same identity proves the lower bound in \eqref{sine}. Applying variation of constants to $y''+q_0y=-(x_0^2-x^2)y$ proves the upper bound; its sine kernel is nonnegative on this interval.

For $0\leq u\leq\sqrt2$, set $M(u)=(1-u^2/E^2)\sin^2(cu)$. It is nondecreasing: half its derivative equals
\[
 \sin(cu)\left[c(1-u^2/E^2)\cos(cu)-\frac u{E^2}\sin(cu)\right]\geq0.
\]
Indeed $c>9/10$, $1-u^2/E^2\geq47/49$, $\cos(cu)\geq\cos\sqrt2\geq7/45$, $0\leq\sin(cu)\leq1$ and $u/E^2<3/98$. The cosine bound is its alternating Taylor lower bound through degree six. The sine upper bound through degree five gives $\sin\sqrt2\leq(7/10)\sqrt2$. Consequently
\[
 M(u)\leq c^2\sin^2(c\sqrt2)\leq\frac{49}{50}c^2.
\]
Using $u=\sqrt E x$, \eqref{sine} and \eqref{initial}, we obtain
\begin{equation}\label{small}
 (E-x^2)y(x)^2\leq\frac{49}{50}A^2\leq\frac{99}{100}EC_n^2
 \quad(0\leq x\leq x_0).
\end{equation}
The lower sine comparison and \eqref{energy} give
\[
 B(0)-B(x_0)\geq\frac{2A^2}{E^2}\int_0^{\sqrt2}u\sin^2u\,du
 >\frac{A^2}{2E^2}.
\]
For the strict last inequality, restrict the integral to $[1,\sqrt2]$, where $\sin u\geq\sin1\geq5/6$; its lower bound is $25/72>1/4$. Thus \eqref{initial} implies $B(x_0)<EC_n^2$. Monotonicity of $B$ proves the bulk bound for $x_0\leq x<\sqrt E$, completing all odd degrees $n\geq3$.

For $n=1$, put $v=x^2\in[0,3]$. Its squared bulk bound is
$\pi v(1-v/3)\leq e^v$. Use $\pi<22/7$, $e^v\geq1+v+v^2/2$ and the positive quadratic
\[
 1+v+v^2/2-(22/7)v(1-v/3)=1-15v/7+65v^2/42>0.
\]
We have therefore proved $|y(x)|\leq C_ng_n(x)$ throughout $0\leq x<\sqrt E$ for every $n\geq1$.

\section*{The clamped tail}
Undo the harmless sign change. For $n\geq3$, we prove $|y(x)|\leq|y(a_n)|$ for $x\geq a_n$. All roots $r_i$ of $p=H_n$ are real, counted with multiplicity. For completeness, if $v_k=H_k(z)$ and $J_k=\operatorname{Im}(v_{k+1}\overline{v_k})$, the physical three-term recurrence gives
\[
 J_0=2\operatorname{Im}z,\qquad
 J_{k+1}=2\operatorname{Im}z\,|v_{k+1}|^2+2(k+1)J_k.
\]
For $\operatorname{Im}z\ne0$, induction gives $(\operatorname{Im}z)J_k>0$ and rules out $v_{k+1}=0$.

Suppose $p(x)\ne0$ and $y'(x)=0$, and put $t=p'/p$, $q=E-x^2$. The Riccati identity $t'=2xt-2n-t^2$, differentiated twice, gives at this point
\[
 t=x,\quad t'=1-q,\quad t''=2x,\quad t'''=2-2q^2.
\]
Since $t=\sum_i(x-r_i)^{-1}$ locally away from the roots, writing $b_i=(x-r_i)^{-1}$ yields
\[
 \sum_i b_i^2=q-1,\qquad \sum_i b_i^3=x,\qquad
 \sum_i b_i^4=(q^2-1)/3.
\]
Cauchy--Schwarz applied to $(b_i)$ and $(b_i^2)$ gives
\begin{equation}\label{peak}
 3E\leq F(q),\qquad F(q)=q^3-q^2+2q+1.
\end{equation}
This uses neither simple roots nor a location beyond all the roots.

The polynomial $F$ is increasing on $\mathbb R$, because
$F'(q)=3(q-1/3)^2+5/3>0$. Let $b=22/7$ and $R(E)=b-b^2/(4E)$. For $E\geq7$,
$\pi-\pi^2/(4E)\leq R(E)$, since their difference is
$(b-\pi)(4E-b-\pi)/(4E)\geq0$. The exact identity
\begin{align*}
 &(3E-F(R(E)))117649E^3\\
 &\quad=352947(E-7)^4+6535179(E-7)^3+40836159(E-7)^2\\
 &\qquad\quad+89231464(E-7)+19560668>0
\end{align*}
contradicts \eqref{peak} if $x\geq a_n$, since $q\leq E-a_n^2=\pi-\pi^2/(4E)$. Thus no nonzero stationary value lies at or beyond $a_n$.

Now $y(x)^2\to0$ as $x\to\infty$. If $y(x)^2>y(a_n)^2$ for some $x>a_n$, choose $R>x$ with $y(R)^2<y(x)^2$. A maximum of $y^2$ on $[a_n,R]$ is positive and interior. Fermat's theorem makes $y'=0$ there, a contradiction. The bulk estimate at $a_n$ supplies the claimed tail.

The remaining degrees have elementary global square bounds
\[
 (2xe^{-x^2/2})^2\leq3/2,\qquad ((4x^2-2)e^{-x^2/2})^2\leq25/4.
\]
The first follows from $ve^{-v}\leq e^{-1}$ and $e\geq8/3$. For the second, calculus gives the possible absolute maxima $2$ and $8e^{-5/4}$; the Taylor bound $e^{5/2}\geq4169/384>256/25$ suffices. These fit the source constants: for $Q_n=1-a_n^2/(2n+1)$, $Q_1=\pi/3-\pi^2/36\leq\pi/4$ and $C_1^2\pi=4$, so $(3/2)Q_1<C_1^2$ using $3<\pi<22/7$. Similarly $Q_2\leq\pi/5$, $C_2=2$ and $\pi<16/5$ give $(25/4)Q_2<C_2^2$.

Combining the bulk and tail estimates and multiplying by $e^{-x^2/2}$ proves Theorem~\ref{main}.

\section*{Attribution and formal certificate}
The source question and normalization are due to Guo and Liu. Classical Hermite identities, real-rootedness, oscillator comparison, Wallis limits, Taylor inequalities and Cauchy--Schwarz are ingredients, not separately claimed results. Foster--Krasikov \cite{FK} gives a different explicit oscillatory-region envelope. We make no assertion that every possible combination of earlier estimates has been excluded, and no worldwide-priority claim.

The Lean certificate has 162 declarations in 17 modules, checked with warnings treated as errors in Lean 4.34.1 and mathlib revision\newline
\texttt{d13f23b723b8a846827a245b89c10fc7d3f11612}. The endpoint \texttt{source\_conjecture\_346243} is the literal source statement, using Rodrigues definitions and the stated Gamma coefficient and cutoff. Initial values, the ODE comparison, strict energy loss, real-root moment obstruction and clamped tail are proved. Its only axioms are \texttt{propext}, \texttt{Classical.choice} and \texttt{Quot.sound}. The endpoint has no residual odd-degree, comparison or energy-loss premise. Independent final review, the extracted packet replay and discovery acceptance are pending.

\small
\begin{thebibliography}{9}
\bibitem{GL} J.~Guo and Y.~Liu, \emph{A modified Euler--Maclaurin formula in 1D and 2D with applications in statistical physics}, Communications in Theoretical Physics \textbf{73} (2021), 075002; arXiv:2004.10441v1. \href{https://doi.org/10.1088/1572-9494/abf1dd}{doi:10.1088/1572-9494/abf1dd}.
\bibitem{FK} W.~H. Foster and I.~Krasikov, \emph{Explicit bounds for Hermite polynomials in the oscillatory region}, LMS Journal of Computation and Mathematics \textbf{3} (2000), 307--314. \href{https://doi.org/10.1112/S1461157000000310}{doi:10.1112/S1461157000000310}.
\end{thebibliography}
\end{document}
